\bibitem{NGFaltings}
G. Faltings, \emph{Endlichkeitss\"atze f\"ur abelsche Variet\"aten
\"uber Zahlk\"orpern}, Invent. Math. \textbf{73} (1983), 349--366,
\S5, S\"atze 3--4, printed p.~360 (PDF p.~12).
\url{https://math.uchicago.edu/~drinfeld/Deligne's_conjecture_Manin_conf/Faltings_argument/Faltings.pdf}.

\bibitem{NGDeligneMilne}
P. Deligne and J. S. Milne, \emph{Tannakian Categories}, in
\emph{Hodge Cycles, Motives, and Shimura Varieties}, LNM 900,
Springer (1982), 101--228.  Corrected transcription of 15 August 2012,
with the original statement numbering: Propositions 6.5, 6.13--6.14,
6.21, 6.23; Theorem 6.25; Proposition 6.26 and Corollary 6.27.
\url{https://www.jmilne.org/math/xnotes/tc2012.pdf}.

\bibitem{NGDeligne82}
P. Deligne (notes by J. S. Milne), \emph{Hodge cycles on abelian
varieties}, in \emph{Hodge Cycles, Motives, and Shimura Varieties},
LNM 900, Springer (1982), 9--100, Proposition 2.9 and Theorem 2.11.
The updated transcription dated 1 October 2018 places these on
pp.~18--19.
\url{https://jmilne.org/math/Documents/Deligne82.pdf}.

\bibitem{NGCommelin}
J. Commelin, \emph{The Mumford--Tate conjecture for products of
abelian varieties}, Algebraic Geometry \textbf{6}(6) (2019), 650--677,
Lemma 3.8 (p.~655); Theorem 4.13 and Remark 4.14 (pp.~657--658);
\S4.16 and Lemma 4.17 (p.~658).
\url{https://content.algebraicgeometry.nl/2019-6/2019-6-028.pdf}.

\bibitem{NGMoonen}
B. Moonen, \emph{Families of Motives and the Mumford--Tate Conjecture},
Example 2.1.13 (p.~10 in the author's 52-page manuscript),
for the Mumford--Tate group of a non-CM elliptic curve.
\url{https://www.math.ru.nl/~bmoonen/Papers/FMMTC.pdf}.

\bibitem{NGNewtonThorne}
J. Newton and J. A. Thorne, \emph{Symmetric power functoriality for
holomorphic modular forms, II}, Publ. Math. IH\'ES \textbf{134}
(2021), 117--152, Theorem A (p.~117), Corollary B (p.~118).
\url{https://www.numdam.org/item/10.1007/s10240-021-00126-4.pdf}.

\bibitem{NGZarhin}
Y. G. Zarhin, \emph{The endomorphism rings of jacobians of cyclic
covers of the projective line}, arXiv:math/0103203v2 (8 September
2002), Theorem 1.1 and Examples 1.3(1), PDF p.~2; genus formula in
\S1, PDF p.~1, and differential description in Remark 3.1,
PDF pp.~7--8.
\url{https://arxiv.org/pdf/math/0103203}.
